\begin{abstract}
% STE descriptive
Mechanistic accounts of modular arithmetic in transformers largely stop where global
invertibility stops. \citet{chen2026multiplication} extend the character account to
$a \cdot b \bmod n$. They treat $\mathbb{Z}/n\mathbb{Z}$ as a monoid, and its Green's
$\mathcal{J}$-classes stratify the input space. Their evidence covers four square-free moduli,
and they describe it as only correlational. Non-square-free moduli have non-regular classes, and
these classes contain nilpotents. \citet{chen2026multiplication} name these moduli as the case
that breaks the reduction to local group characters.

We show that the reduction does not break. With $x = d \cdot u$, a $\mathcal{J}$-class is a
torsor under the \emph{local} units $(\mathbb{Z}/(n/d)\mathbb{Z})^{\times}$. Under the global
units the action is transitive but not free. So every stratum is multiplication in a smaller
ring, carried on that stratum's own local group $G_{m} = (\mathbb{Z}/(n/m)\mathbb{Z})^{*}$.
Nilpotents collapse a stratum onto a smaller local group. They do not destroy its character structure.

We measure this across $23$ moduli, $14$ of them non-square-free and $6$ of them prime powers.
A transformer runs the same character-based clock on every stratum. The measurement covers $32$
strata, $157$ stratum-run measurements and $12$ distinct local groups. We committed its
criteria before any per-stratum number existed. At $n = 165$ the strata raise the share of the
multiplication table that the mechanism accounts for from $23.5\%$ to $82.3\%$.

It is also the \emph{same} clock, and not one clock per stratum. Strata that share a local group
select exactly the same characters of it in $158$ of $160$ pre-registered pairs. Against this, a
$10{,}000$-draw null is at its floor, and a control that destroys the structure reads $0$ of
$160$.

The output depends on those frequencies, and they are not incidental. A $10{,}000$-draw
permutation ablation of the logit spectrum puts the key set beyond \emph{every} draw at $133$ of
$157$ stratum-measurements. It gives $\hat{p} \le 0.0023$ at all $157$, with excluded/baseline
from $3.98 \times 10^{2}$ to $2.71 \times 10^{8}$. It also holds in $28$ of $29$ further
single-modulus runs, and we name the one exception. It reaches $2^{7}$, the deepest nilpotency
tested, in product-character coordinates where no discrete log exists.

This is an intervention on the saved logit spectrum, the restricted/excluded protocol of
\citet{nanda2023progress}. It re-runs no forward pass, ablates no weight and patches no
activation. So it establishes what the output depends on, and not which internal component
computes it. We therefore claim no more than that protocol earns, and we name it wherever we
cite the evidence.

At two moduli we do claim more: we delete those characters from the \emph{embedding} and run
the network forward unmodified. This takes unit-grid accuracy to $0.0089$--$0.0091$ at the prime
$n = 113$ (chance $1/112$). It takes it to $0.0000$--$0.0120$ at $n = 121 = 11^{2}$ (chance
$1/110$), and the edit leaves the non-regular stratum of that modulus intact. Both results hold
on three of three seeds at each modulus. The complementary edit keeps only those characters and
leaves accuracy at $0.9596$ or above. This is a weight-space intervention with both
nonlinearities downstream of the edit, on checkpoints that reproduce the archived runs
bit-identically.

The algebra of $n$ changes how many strata there are and how they nest, and not the mechanism
that runs on them. Two negatives bound the claim. The criterion that our own pre-registration
named primary does not discriminate a grokked model from a failed one. So the verdict rests on
the three criteria that do. The clearest early signal that we found also rises identically in
runs that never generalise. It is therefore not a progress measure.
% STE end
\end{abstract}
